\subsection{RQ1.5b: AKS Multi-Node Sensitivity}
\label{sec:analysis-rq15b}

The results in \cref{sec:rq15b} support three bounded findings for the AKS
multi-node sensitivity stage.

\paragraph{Finding 1: Multi-node placement preserves the AKS condition ordering.}
With every chain segment forced onto a distinct node and the benchmark client
on a dedicated node, the complete RQ1.5 ordering is preserved
(\texttt{monolithic\_k8s\_1svc}~$<$~\texttt{chain\_2svc}~$<$~\texttt{chain\_3svc}~$<$~\texttt{chain\_4svc}~$<$~\texttt{chain\_5svc}).
The bounded-non-linearity pattern from RQ1.4 and RQ1.5 also reproduces under
multi-node execution: the
\texttt{chain\_4svc}~$\rightarrow$~\texttt{chain\_5svc} increment is 2.502~ms,
broadly comparable to the preceding 4.562~ms step
(\cref{tab:rq15b_results}). The architectural ordering is therefore not an
artefact of single-node colocation.

\paragraph{Finding 2: The inter-node penalty is positive and bounded.}
\Cref{tab:rq15b_delta} compares the multi-node run condition-by-condition with
the frozen single-node RQ1.5 baseline. The monolithic condition gains 2.305~ms
of latency relative to the single-node monolithic. This is compatible with
moving the 588~KiB input tensor across the Azure VNet rather than across a
loopback socket. The chained conditions add per-condition deltas between
2.131~ms (\texttt{chain\_3svc}) and 4.086~ms (\texttt{chain\_5svc}). The
multi-node penalty is therefore a measured bound rather than an open
acknowledgement: the worst-case single-node-to-multi-node delta on the
evaluated chain family is 4.086~ms (\texttt{chain\_5svc}), and the monolithic
baseline alone gains 2.305~ms.

\paragraph{Finding 3: Round-to-round consistency remains low.}
Round CVs for the chained conditions in RQ1.5b are 0.0036--0.0127, comparable
to and in some cases tighter than the single-node RQ1.5 chained CVs of
0.0019--0.0061 (\cref{tab:rq15b_rounds,tab:rq15_rounds}). The monolithic
condition shows a similar profile (0.0084 vs 0.0117). Cross-round stability is
therefore preserved when the topology is distributed across nodes.

\paragraph{Answer to RQ1.5b.}
Under enforced multi-node placement, the same five-condition chain family
preserves the monotonic overhead progression seen in RQ1.4 and RQ1.5:
per-condition mean latency increases from 73.108~ms (monolithic) to 84.605~ms
(\texttt{chain\_5svc}), and the
\texttt{chain\_4svc}~$\rightarrow$~\texttt{chain\_5svc} step is +2.502~ms.
The per-condition multi-node penalty relative to the single-node RQ1.5
baseline ranges from +2.131~ms to +4.086~ms across the chained conditions and
+2.305~ms for the monolithic baseline. The single-node RQ1.4 / RQ1.5 design
therefore underestimates true production latency by a bounded margin that is
now quantified per condition rather than acknowledged abstractly.
